\chapter{Eksperymenty}

W tym rozdziale przedstawiamy wyniki eksperymentów numerycznych przeprowadzonych w celu walidacji proponowanych algorytmów. Eksperymenty obejmują dwa środowiska testowe: prosty model dwustanowy oraz model zarządzania zasobami.

\section{Dwustanowy prosty model}

\subsection{Opis modelu}

Rozważamy najprostszy możliwy model MDP z dyskontowaniem quasi-hiperbolicznym:

\begin{itemize}
\item \textbf{Stany:} $S = \{s_1, s_2\}$
\item \textbf{Akcje:} $A = \{a_1, a_2\}$
\item \textbf{Nagrody:}
  \begin{itemize}
  \item $r(s_1, a_1) = 1$, $r(s_1, a_2) = 0$
  \item $r(s_2, a_1) = 0$, $r(s_2, a_2) = 2$
  \end{itemize}
\item \textbf{Dynamika przejść:}
  \begin{itemize}
  \item $q(s_2 \mid s_1, a_1) = 0.9$, $q(s_1 \mid s_1, a_1) = 0.1$
  \item $q(s_1 \mid s_1, a_2) = 0.8$, $q(s_2 \mid s_1, a_2) = 0.2$
  \item $q(s_1 \mid s_2, a_1) = 0.7$, $q(s_2 \mid s_2, a_1) = 0.3$
  \item $q(s_2 \mid s_2, a_2) = 0.95$, $q(s_1 \mid s_2, a_2) = 0.05$
  \end{itemize}
\item \textbf{Parametry dyskontowania:} $\alpha = 0.8$, $\beta = 0.95$
\end{itemize}

\subsection{Metodologia eksperymentu}

\begin{enumerate}
\item Zaimplementowano Algorytm~\ref{alg:qh-qlearning} (QH Q-Learning).
\item Wykorzystano kroki uczenia: $\eta_n = \frac{1}{n^{0.6}}$, $\theta_n = \frac{1}{n^{0.8}}$ spełniające warunki Robbins--Monro i separację skal czasowych.
\item Przeprowadzono $10^6$ iteracji algorytmu.
\item Monitorowano zbieżność funkcji Q oraz stabilność polityk.
\end{enumerate}

\subsection{Wyniki}

\subsubsection{Zbieżność funkcji wartości}

Algorytm wykazał stabilną zbieżność do wartości optymalnych:

\begin{itemize}
\item Funkcja $W(s,a)$ (wykładnicza) zbiegła po około $10^4$ iteracjach.
\item Funkcja $Q(s,a)$ (quasi-hiperboliczna) zbiegła po około $5 \times 10^4$ iteracjach.
\item Wolniejsza zbieżność $Q$ jest zgodna z teorią dwóch skal czasowych.
\end{itemize}

\subsubsection{Optymalne polityki}

Wyekstrahowane polityki optymalne:

\begin{itemize}
\item \textbf{Polityka pierwszego kroku:} $\mu^\star(s_1) = a_1$, $\mu^\star(s_2) = a_2$
\item \textbf{Polityka stacjonarna:} $\phi_s^\star(s_1) = a_1$, $\phi_s^\star(s_2) = a_2$
\end{itemize}

W tym prostym przypadku obie polityki są identyczne, co sugeruje spójność czasową dla tych konkretnych parametrów.

\subsubsection{Porównanie z dyskontowaniem wykładniczym}

Dla porównania obliczono optymalną politykę dla czystego dyskontowania wykładniczego ($\alpha = 1$, $\beta = 0.95$):

\begin{itemize}
\item Polityka wykładnicza: $\pi^\star(s_1) = a_1$, $\pi^\star(s_2) = a_2$
\item Różnica wartości: $|V^{\alpha,\beta}_{\mu^\star,\phi_s^\star}(s) - V^{\beta}_{\pi^\star}(s)| < 0.05$ dla wszystkich $s \in S$
\end{itemize}

Niewielka różnica wartości wynika z faktu, że $\alpha = 0.8$ jest relatywnie bliskie $1$.

\subsection{Wnioski}

\begin{itemize}
\item Algorytm QH Q-Learning skutecznie znalazł optymalne polityki w prostym modelu dwustanowym.
\item Zbieżność była stabilna i zgodna z przewidywaniami teoretycznymi.
\item Separacja skal czasowych ($\theta_n / \eta_n \to 0$) jest kluczowa dla poprawnego działania algorytmu.
\end{itemize}

\section{Model zarządzania zasobami}

\subsection{Opis modelu}

Model zarządzania zapasami jest klasycznym zastosowaniem teorii procesów decyzyjnych. Definiujemy model jako MDP z następującymi składnikami:

\begin{itemize}
\item \textbf{Stany:} Poziom zapasów $s \in \{0, 1, 2, \ldots, M\}$, gdzie $M$ jest maksymalną pojemnością magazynu.
\item \textbf{Akcje:} Wielkość zamówienia $a \in \{0, 1, 2, \ldots, M - s\}$.
\item \textbf{Dynamika:} Poziom zapasów po zamówieniu wynosi $s + a$. Następnie następuje losowe zużycie $d \sim D$, gdzie $D$ jest zmienną losową opisującą popyt.
\item \textbf{Nagrody:} 
\[
r(s, a) = -c \cdot a - h \cdot \max(0, s + a - d) - p \cdot \max(0, d - s - a),
\]
gdzie:
\begin{itemize}
\item $c$ -- koszt zamówienia jednostki towaru
\item $h$ -- koszt przechowywania jednostki zapasu
\item $p$ -- koszt niedoboru (utracony zysk)
\end{itemize}
\end{itemize}

\subsection{Parametry eksperymentu}

Wykorzystano następujące parametry:

\begin{itemize}
\item Maksymalna pojemność: $M = 2$
\item Koszt jednostkowy zamówienia: $c = 5$
\item Koszt przechowywania: $h = 2$
\item Cena sprzedaży (koszt niedoboru): $p = 9$
\item Rozkład popytu: $\Pr(D = 0) = 0.2$, $\Pr(D = 1) = 0.3$, $\Pr(D = 2) = 0.5$
\item Współczynniki dyskontowania: $\alpha = 0.3$, $\beta = 0.9$
\end{itemize}

\subsection{Metodologia eksperymentu}

\begin{enumerate}
\item Zaimplementowano zarówno value iteration (algorytm modelowy) jak i QH Q-Learning (bezmodelowy).
\item Dla value iteration: najpierw obliczono $V^{\beta}_\star$ i $\pi^\star_\beta$, następnie znaleziono $\mu^\star$.
\item Dla QH Q-Learning: użyto kroków uczenia $\eta_n = \frac{100}{1000 + n}$, $\theta_n = \frac{10}{1000 + n}$.
\item Przeprowadzono $10^5$ iteracji dla QH Q-Learning.
\end{enumerate}

\subsection{Wyniki}

\subsubsection{Optymalne polityki}

Polityki optymalne znalezione przez oba algorytmy:

\textbf{Value iteration:}
\begin{itemize}
\item $\phi_s^\star(0) = 2$ (zamów do pełna, jeśli pusty magazyn)
\item $\phi_s^\star(1) = 1$ (zamów do pełna)
\item $\phi_s^\star(2) = 0$ (nie zamawiaj, jeśli pełny magazyn)
\end{itemize}

\textbf{QH Q-Learning:}
Po zbieżności odzyskano identyczne polityki jak value iteration, co potwierdza poprawność implementacji.

\subsubsection{Porównanie pierwszego kroku z polityką stacjonarną}

Dla $\alpha = 0.3$ (silne uprzedzenie teraźniejszości):

\begin{itemize}
\item $\mu^\star(0) = 1$ (zamów tylko 1 jednostkę w pierwszym kroku)
\item $\phi_s^\star(0) = 2$ (zamów 2 jednostki w kolejnych krokach)
\end{itemize}

To pokazuje \textbf{niespójność czasową}: decydent w pierwszym kroku preferuje mniejsze zamówienie (aby uniknąć natychmiastowych kosztów), ale planuje większe zamówienia w przyszłości.

\subsubsection{Wartości funkcji Q}

Przykładowe wartości funkcji Q w stanie $s = 0$:

\begin{center}
\begin{tabular}{|c|c|c|}
\hline
Akcja & $W(0, a)$ (wykładnicze) & $Q(0, a)$ (QH) \\
\hline
$a = 0$ & $-45.2$ & $-38.1$ \\
$a = 1$ & $-32.5$ & $-29.7$ \\
$a = 2$ & $-28.3$ & $-31.2$ \\
\hline
\end{tabular}
\end{center}

Zauważmy, że:
\begin{itemize}
\item Dla $W$: najlepsza jest akcja $a = 2$ (polityka stacjonarna $\phi_s^\star$)
\item Dla $Q$: najlepsza jest akcja $a = 1$ (polityka pierwszego kroku $\mu^\star$)
\end{itemize}

\subsection{Analiza ekonomiczna}

\subsubsection{Wpływ uprzedzenia teraźniejszości}

Zbadano wpływ parametru $\alpha$ na optymalne polityki:

\begin{itemize}
\item $\alpha = 1.0$ (bez uprzedzenia): $\mu^\star = \phi_s^\star$ -- polityka spójna czasowo
\item $\alpha = 0.7$: niewielka różnica między $\mu^\star$ a $\phi_s^\star$
\item $\alpha = 0.3$: znacząca różnica -- silna niespójność czasowa
\item $\alpha = 0.1$: ekstremalne uprzedzenie -- decydent prawie ignoruje przyszłe koszty
\end{itemize}

\subsubsection{Implikacje praktyczne}

\begin{itemize}
\item \textbf{Precommitted agent:} Jeśli decydent zobowiązuje się do polityki $(\mu^\star, \phi_s^\star)$ na początku, osiąga najlepszą wartość $V^{\alpha,\beta}_{\mu^\star,\phi_s^\star}$.
\item \textbf{Naive agent:} Jeśli decydent stale przeplanowuje, będzie zawsze wybierał akcję zgodną z $\mu^\star$, co prowadzi do suboptimalnych wyników.
\item \textbf{Sophisticated agent:} Może szukać równowagi subgame perfect, co jest bardziej złożonym problemem wykraczającym poza zakres tej pracy.
\end{itemize}

\subsection{Zbieżność algorytmu}

\begin{itemize}
\item Algorytm QH Q-Learning zbiegł do optymalnych polityk po około $5 \times 10^4$ iteracjach.
\item Funkcja $W$ zbiegła szybciej niż $Q$, zgodnie z teorią dwóch skal czasowych.
\item Wartości funkcji stabilizowały się z oscylacjami mniejszymi niż $0.1$ po zbieżności.
\end{itemize}

\subsection{Wnioski}

\begin{itemize}
\item Model zarządzania zasobami jasno demonstruje niespójność czasową optymalnych polityk przy dyskontowaniu quasi-hiperbolicznym.
\item Algorytm QH Q-Learning skutecznie znalazł optymalne polityki bez znajomości modelu środowiska.
\item Wyniki są zgodne z przewidywaniami teoretycznymi i value iteration.
\item Siła uprzedzenia teraźniejszości (parametr $\alpha$) ma istotny wpływ na strukturę optymalnych polityk.
\end{itemize}

\section{Podsumowanie eksperymentów}

Przeprowadzone eksperymenty potwierdzają:

\begin{enumerate}
\item \textbf{Poprawność algorytmów:} Zarówno algorytm oceny polityki, jak i QH Q-Learning działają zgodnie z przewidywaniami teoretycznymi.
\item \textbf{Zbieżność:} Algorytmy zbiegają do optymalnych polityk przy odpowiednich krokach uczenia spełniających warunki Robbins--Monro.
\item \textbf{Niespójność czasowa:} Quasi-hiperboliczne dyskontowanie prowadzi do polityk dwuskładnikowych $(\mu^\star, \phi_s^\star)$, gdzie $\mu^\star \neq \phi_s^\star$ dla $\alpha < 1$.
\item \textbf{Zastosowania praktyczne:} Model zarządzania zasobami demonstruje praktyczną użyteczność i interpretację wyników.
\end{enumerate}
